% ot_kl_divergence.tex
% "Couplings, Wasserstein Distances, and Kullback-Leibler Divergence"
% Thomas J. Sargent
% April 2026

\documentclass[12pt]{article}

\usepackage{amsmath,amssymb}
\usepackage{booktabs}
\usepackage{array}
\usepackage[margin=1.2in]{geometry}
\usepackage{setspace}
\usepackage{microtype}
\usepackage{hyperref}
\usepackage{natbib}
\usepackage{parskip}

\onehalfspacing

\hypersetup{
  colorlinks = true,
  linkcolor  = blue,
  citecolor  = blue,
  urlcolor   = blue
}

% Convenient shorthand
\newcommand{\KL}[2]{D_{\mathrm{KL}}\!\left(#1\,\|\,#2\right)}
\newcommand{\TV}[2]{\left\|#1 - #2\right\|_{\mathrm{TV}}}
\newcommand{\Wass}[3]{W_{#1}\!\left(#2,\,#3\right)}
\newcommand{\E}[2]{\mathbb{E}_{#1}\!\left[#2\right]}
\newcommand{\calP}{\mathcal{P}}
\newcommand{\calH}{\mathcal{H}}
\newcommand{\R}{\mathbb{R}}
\newcommand{\ind}{\mathbf{1}}

\title{Couplings, Wasserstein Distances, and\\
       Kullback-Leibler Divergence}
\author{Thomas J.\ Sargent}
\date{April 2026}

\begin{document}
\maketitle

\begin{abstract}
We study two prominent ways of measuring discrepancy between probability
distributions---the Kullback-Leibler (KL) divergence and the Wasserstein
distance---on a shared compact metric space $X$.  After defining each
object precisely and cataloguing its key properties, we examine how they
differ in geometric sensitivity, in their requirements on absolute
continuity, and in their dual variational representations.  A central
question is whether the KL divergence, like the Wasserstein distance,
can be understood as an optimisation over \emph{couplings} (joint
distributions with prescribed marginals).  We argue that it cannot in
general---KL uses a fixed pointwise comparison at the same location
$x \in X$ and involves no transport of mass---but that the two objects
are linked through two related constructions:
\emph{entropy-regularised optimal transport} (in which KL appears as
a penalty on couplings) and the classical \emph{Schr\"odinger bridge}
problem (in which KL on joint distributions replaces the transport cost).
We close with inequalities that bound one quantity in terms of the other,
including Pinsker's inequality and Talagrand's transport inequality.
\end{abstract}

\tableofcontents
\newpage

%======================================================================
\section{Introduction}

Let $X$ be a compact metric space equipped with its Borel
$\sigma$-algebra $\mathcal{B}(X)$ and a reference measure
$\lambda$ (typically Lebesgue measure on a subset of $\R^d$,
or the counting measure on a finite set).  Denote by
$\calP(X)$ the set of Borel probability measures on $X$.

We consider two probability density functions $f$ and $g$ with
respect to $\lambda$, so that $f, g \geq 0$ and
$\int_X f\,d\lambda = \int_X g\,d\lambda = 1$.  We say that
$f$ and $g$ are \emph{mutually absolutely continuous} if
$f(x) = 0 \Leftrightarrow g(x) = 0$ $\lambda$-almost
everywhere, written $f \ll g$ and $g \ll f$.  In the language
of measures, letting $\mu, \nu \in \calP(X)$ denote the
measures with densities $f$ and $g$ respectively, mutual
absolute continuity means $\mu \ll \nu$ and $\nu \ll \mu$.

Two fundamental and ubiquitous ways to measure how far $f$ is
from $g$ are:
\begin{enumerate}
  \item The \textbf{Kullback-Leibler divergence}
    $\KL{f}{g}$, rooted in information theory and statistics.
  \item The \textbf{Wasserstein distance} $\Wass{p}{f}{g}$,
    rooted in optimal transport theory.
\end{enumerate}
Both reduce to zero if and only if $f = g$ almost everywhere.
Beyond this they are, in many ways, measuring different things:
KL divergence is insensitive to the ground metric $d$ on $X$
(the metric on the \emph{sample space}); Wasserstein is
exquisitely sensitive to it.  On the other hand, KL divergence
does induce a rich geometry on the \emph{space of probability
distributions} $\calP(X)$ via the Fisher information metric---the
subject of Amari's information geometry \citep{Amari2016,AmariNagaoka2000}---a
distinct geometry from that generated by $W_2$.  KL divergence requires absolute
continuity; Wasserstein does not.  KL divergence is defined
directly in terms of the two densities evaluated at the same
point; Wasserstein is defined as an infimum over all joint
distributions (\emph{couplings}) of a pair having $f$ and $g$
as their marginals.

Understanding these similarities and differences is essential
for applications in statistics, machine learning, information
theory, mathematical physics, and economics.  This note aims to
give a self-contained exposition.

%======================================================================
\section{Kullback-Leibler Divergence}
\label{sec:kl}

\subsection{Definition}

\begin{equation}\label{eq:kl_def}
  \boxed{\KL{f}{g} = \int_X f(x)\log\frac{f(x)}{g(x)}\,d\lambda(x).}
\end{equation}
Here the integrand is defined to be $0$ when $f(x) = 0$ (by
$0\log 0 = 0$), and $+\infty$ when $f(x) > 0 = g(x)$.  Thus
$\KL{f}{g} \in [0, +\infty]$, and it is finite precisely when
$f \ll g$.  Under mutual absolute continuity, it is always
finite.

The KL divergence can be decomposed as
\begin{equation}\label{eq:kl_entropy}
  \KL{f}{g} = -H(f) + H_{\mathrm{cross}}(f,g),
\end{equation}
where $H(f) = -\int f\log f\,d\lambda$ is the
\emph{differential entropy} of $f$, and
$H_{\mathrm{cross}}(f,g) = -\int f\log g\,d\lambda$ is the
\emph{cross-entropy} of $f$ relative to $g$.

\subsection{Properties}

\paragraph{Non-negativity (Gibbs' inequality).}
$\KL{f}{g} \geq 0$, with equality if and only if $f = g$
$\lambda$-a.e.  This follows from Jensen's inequality applied to
the strictly convex function $\phi(t) = t\log t$:
\begin{equation}\label{eq:gibbs}
  \KL{f}{g} = \int_X g(x)\,\phi\!\left(\frac{f(x)}{g(x)}\right)
  d\lambda(x) \geq \phi\!\left(\int_X g(x)
  \frac{f(x)}{g(x)}\,d\lambda(x)\right) = \phi(1) = 0.
\end{equation}

\paragraph{Asymmetry.}
In general $\KL{f}{g} \neq \KL{g}{f}$.  The
\emph{symmetrised KL divergence} (also called the $J$-divergence)
is
\begin{equation}\label{eq:j_div}
  J(f,g) = \KL{f}{g} + \KL{g}{f}
  = \int_X \bigl(f(x) - g(x)\bigr)\log\frac{f(x)}{g(x)}\,d\lambda(x),
\end{equation}
which is symmetric and non-negative.

\paragraph{Not a metric.}
$\KL{f}{g}$ is not a metric: it is not symmetric, and it does
not satisfy the triangle inequality.

\paragraph{Convexity.}
The map $(f,g) \mapsto \KL{f}{g}$ is jointly convex: for any
$\alpha \in [0,1]$,
\[
  \KL{\alpha f_1 + (1-\alpha)f_2}{\alpha g_1 + (1-\alpha)g_2}
  \leq \alpha\,\KL{f_1}{g_1} + (1-\alpha)\,\KL{f_2}{g_2}.
\]

\paragraph{Data processing inequality.}
Let $T: X \to Y$ be any measurable map and let $T_\#\mu$,
$T_\#\nu$ denote the push-forward measures.  Then
\begin{equation}\label{eq:dpi}
  \KL{T_\#\mu}{T_\#\nu} \leq \KL{\mu}{\nu}.
\end{equation}
Equivalently, processing the data reduces (or preserves)
information.

\paragraph{Chain rule.}
Let $\pi_f$ and $\pi_g$ be joint distributions over $X \times X$
with $x$-marginals $f$ and $g$ respectively.  Denote by
$f(\cdot|x)$ and $g(\cdot|x)$ their conditional distributions given
the first coordinate.  Then
\begin{equation}\label{eq:chainrule}
  \KL{\pi_f}{\pi_g}
  = \KL{f_1}{g_1}
  + \int_X f_1(x)\,\KL{f(\cdot|x)}{g(\cdot|x)}\,d\lambda(x),
\end{equation}
where $f_1$ and $g_1$ are the first marginal densities.  The
chain rule shows that KL divergence is additive over conditional
levels of a hierarchy.

\paragraph{Sensitivity to tails.}
Because the integrand in \eqref{eq:kl_def} is
$f \log(f/g) = f \log f - f \log g$, regions where $g$ is
very small but $f$ is not contribute large (potentially
infinite) quantities.  KL divergence is thus very sensitive to
tail behaviour.  The Wasserstein distance, by contrast, is
insensitive to extreme tails (on a compact space, it is always
finite).

\subsection{Variational Representation}
\label{sec:kl_var}

The \emph{Donsker-Varadhan variational formula}
\citep{DonskerVaradhan1975} expresses KL divergence as the
Legendre transform of the log-moment-generating functional:
\begin{equation}\label{eq:dv}
  \boxed{\KL{f}{g}
  = \sup_{h \in L^1(f)}
  \left[\E{f}{h(X)} - \log\E{g}{e^{h(X)}}\right].}
\end{equation}
The supremum is attained at the function $h^*(x) =
\log\bigl(f(x)/g(x)\bigr)$.  Substituting $h^*$:
\[
  \E{f}{h^*} - \log\E{g}{e^{h^*}}
  = \E{f}{\log\frac{f}{g}} - \log\int g\cdot\frac{f}{g}\,d\lambda
  = \KL{f}{g} - \log 1 = \KL{f}{g}. \checkmark
\]
The formula \eqref{eq:dv} is extremely useful because the right-hand
side depends only on expectations under $f$ and $g$ separately,
with no need to evaluate the ratio $f/g$ directly.  This makes it
numerically tractable and is the basis of many estimators of KL
divergence from samples.

\subsection{The \texorpdfstring{$f$}{f}-Divergence Family}
\label{sec:f_div}

KL divergence is a member of the \emph{$f$-divergence} family
introduced by \citet{Csiszar1975}.  Given a convex function
$\phi: (0,\infty) \to \R$ with $\phi(1) = 0$, the
$\phi$-divergence is
\begin{equation}\label{eq:f_div}
  D_\phi(f \| g)
  = \int_X g(x)\,\phi\!\left(\frac{f(x)}{g(x)}\right)d\lambda(x).
\end{equation}
Special cases are listed in Table~\ref{tab:f_divs}.

\begin{table}[ht]
\centering
\caption{Members of the $f$-divergence family defined by
\eqref{eq:f_div}.  All satisfy $D_\phi \geq 0$ with equality
iff $f = g$ a.e.}
\label{tab:f_divs}
\medskip
\begin{tabular}{p{4.2cm}p{3.6cm}p{4.5cm}}
\toprule
Divergence & Generator $\phi(t)$ & Formula \\
\midrule
KL (forward) & $t\log t$ & $\int f\log(f/g)$ \\
KL (reverse) & $-\log t$ & $\int g\log(g/f)$ \\
Total variation & $\frac{1}{2}|t-1|$ & $\frac{1}{2}\int|f-g|$ \\
$\chi^2$-divergence & $(t-1)^2$ & $\int(f-g)^2/g$ \\
Hellinger squared & $(\sqrt{t}-1)^2$ & $1 - \int\sqrt{fg}$ \\
Jensen--Shannon & $(t+1)\log\frac{2}{t+1}+t\log t$ & $\frac{1}{2}[\KL{f}{m}+\KL{g}{m}]$, $m=\frac{f+g}{2}$ \\
\bottomrule
\end{tabular}
\end{table}

$f$-divergences share several structural features:
non-negativity, joint convexity, and the data-processing
inequality \eqref{eq:dpi}.  Crucially for our comparison, they
are all \emph{reference-measure-based}: they evaluate the ratio
$f(x)/g(x)$ at the same point $x$.  No transport is involved.

%======================================================================
\section{Couplings}
\label{sec:couplings}

\subsection{Marginals and the Set of Couplings}

Let $\mu$ and $\nu$ be two probability measures in $\calP(X)$.
A \emph{coupling} of $\mu$ and $\nu$ is a probability measure
$\pi \in \calP(X \times X)$ whose first marginal is $\mu$ and
whose second marginal is $\nu$:
\begin{equation}\label{eq:coupling_def}
  \pi(A \times X) = \mu(A)
  \quad\text{and}\quad
  \pi(X \times B) = \nu(B)
  \quad \text{for all } A, B \in \mathcal{B}(X).
\end{equation}
The set of all couplings is denoted
\begin{equation}
  \Pi(\mu,\nu) = \bigl\{\pi \in \calP(X\times X) :
    \text{\eqref{eq:coupling_def} holds}\bigr\}.
\end{equation}
$\Pi(\mu,\nu)$ is always non-empty (it contains the
\emph{independent coupling} $\mu \otimes \nu$ defined by
$(\mu\otimes\nu)(A\times B) = \mu(A)\nu(B)$), convex, and
compact in the weak topology when $X$ is compact.

If $\mu$ and $\nu$ both have densities $f$ and $g$ with respect
to $\lambda$, a coupling $\pi \in \Pi(\mu,\nu)$ with
joint density $\pi(x,y)$ satisfies:
\begin{equation}\label{eq:marginal_constraints}
  \int_X \pi(x,y)\,d\lambda(y) = f(x),
  \qquad
  \int_X \pi(x,y)\,d\lambda(x) = g(y).
\end{equation}
These are the \emph{marginal constraints} of the optimal
transport problem.

\subsection{Examples of Couplings}

\paragraph{1.~The independent coupling.}
$\pi_\ind(x,y) = f(x)g(y)$.  Under this coupling, $X$ and $Y$
are independent: knowing $X = x$ conveys no information about
$Y$.  The independent coupling is the unique maximiser of the
entropy $H(\pi) = -\int\pi\log\pi$ subject to the marginal
constraints.

\paragraph{2.~Deterministic (Monge) couplings.}
Given a \emph{transport map} $T: X \to X$ with $T_\#\mu = \nu$
(i.e., $\nu(B) = \mu(T^{-1}(B))$ for all $B$), the
corresponding coupling is
\begin{equation}\label{eq:monge_coupling}
  \pi_T(A \times B) = \mu(A \cap T^{-1}(B)).
\end{equation}
In terms of density: $\pi_T(x,y) = f(x)\,\delta(y - T(x))$.
Under this coupling, $Y = T(X)$ is a deterministic function of
$X$.  The marginal constraint on $y$ requires $T_\#\mu = \nu$,
which is the condition defining a valid transport map.

\paragraph{3.~The diagonal coupling.}
When $\mu = \nu$ (i.e., $f = g$), the \emph{diagonal
coupling} $\pi_\Delta(x,y) = f(x)\delta(y-x)$ has both
marginals equal to $f$.  Under this coupling $X = Y$ almost
surely.  Note that the diagonal coupling is a special case of
a Monge coupling with $T = \mathrm{id}$, and it only belongs
to $\Pi(\mu,\nu)$ when $\mu = \nu$.

\paragraph{4.~An anti-monotone coupling.}
On $X = [0,1]$, pair the $\alpha$-quantile of $f$ with the
$(1-\alpha)$-quantile of $g$ for each $\alpha \in [0,1]$.
This coupling is ``anti-aligned'' and in general incurs larger
transport cost than the identity coupling; it is used as a
worst-case comparison.

\paragraph{5.~The optimal coupling.}
For a given cost function $c: X \times X \to [0,\infty)$, the
coupling that minimises $\int c\,d\pi$ over $\Pi(\mu,\nu)$ is
the \emph{optimal (Kantorovich) coupling}.  Its existence and
characterisation is the central object of optimal transport
theory, to which we now turn.

%======================================================================
\section{Optimal Transport}
\label{sec:ot}

\subsection{The Monge Problem}

Gaspard Monge posed the original mass-transportation problem
in 1781: find a measurable map $T: X \to X$ that transports
$\mu$ to $\nu$ while minimising the total cost:
\begin{equation}\label{eq:monge}
  \inf_{T:\, T_\#\mu = \nu}
  \int_X c(x, T(x))\,d\mu(x).
\end{equation}
Here $c: X \times X \to [0,\infty)$ is a \emph{cost function}
measuring the expense of moving a unit of mass from $x$ to
$T(x)$.  Typical choices are $c(x,y) = d(x,y)$ (transport
distance proportional to displacement) or $c(x,y) = d(x,y)^2$
(quadratic cost, penalising large displacements more severely).

The Monge problem has two difficulties.  First, the feasible set
$\{T: T_\#\mu = \nu\}$ may be empty: for example, if $\mu$ is a
Dirac mass $\delta_{x_0}$ and $\nu$ is not, there is no map
sending $\mu$ to $\nu$.  Second, even when a map exists, the
infimum may not be attained.

\subsection{The Kantorovich Relaxation}

\citet{Kantorovich1942} relaxed the Monge problem by allowing
\emph{mass splitting}: instead of a deterministic map, seek a
coupling $\pi \in \Pi(\mu,\nu)$ that minimises the expected
cost:
\begin{equation}\label{eq:kantorovich}
  \boxed{C(\mu,\nu; c)
  = \inf_{\pi \in \Pi(\mu,\nu)}
  \int_{X\times X} c(x,y)\,d\pi(x,y).}
\end{equation}
The Kantorovich problem is a \emph{linear programme} over the
convex set $\Pi(\mu,\nu)$, and it always has a solution when
$c$ is lower semicontinuous and $X$ is compact.  Every Monge
coupling is a special case (it restricts $\pi$ to the set
$\{(x,T(x)): x \in X\}$), so the Kantorovich infimum is at
most the Monge infimum.  On well-behaved spaces the two
coincide.

\subsection{Kantorovich Duality}

The Kantorovich problem has a beautiful dual formulation.
Define a pair of functions $(\phi, \psi): X \times X \to \R$
to be $c$-\emph{admissible} if
$\phi(x) + \psi(y) \leq c(x,y)$ for all $(x,y)$.
The \emph{Kantorovich dual} is:
\begin{equation}\label{eq:kantorovich_dual}
  C(\mu,\nu; c)
  = \sup_{\phi + \psi \leq c}
  \left[\int_X \phi(x)\,d\mu(x) + \int_X \psi(y)\,d\nu(y)\right].
\end{equation}
Equality between \eqref{eq:kantorovich} and
\eqref{eq:kantorovich_dual} (strong duality) holds under mild
conditions \citep{Villani2009}.  The dual variables $\phi$ and
$\psi$ are called \emph{Kantorovich potentials}; they play the
role of prices in an equilibrium interpretation.

For the quadratic cost $c(x,y) = |x-y|^2$ on $\R^d$, the
$c$-transform $\phi^c(y) = \inf_{x}[|x-y|^2 - \phi(x)]$
yields $\psi = \phi^c$ at the optimum, and both $\phi$ and
$-\phi^c$ are convex functions.  Brenier's theorem
\citep{Brenier1991} then gives: if $\mu$ is absolutely
continuous, the optimal map is $T = \nabla\phi$ for a convex
function $\phi$---a gradient of convex potential.

%======================================================================
\section{Wasserstein Distance}
\label{sec:wasserstein}

\subsection{Definition}

Let $(X, d)$ be a compact metric space and fix $p \geq 1$.
The \emph{Wasserstein-$p$ distance} between $\mu$ and $\nu$
in $\calP(X)$ is
\begin{equation}\label{eq:wasserstein}
  \boxed{W_p(\mu, \nu)
  = \left(\inf_{\pi \in \Pi(\mu,\nu)}
    \int_{X\times X} d(x,y)^p\,d\pi(x,y)\right)^{1/p}.}
\end{equation}
This is the Kantorovich cost \eqref{eq:kantorovich} with
$c(x,y) = d(x,y)^p$, raised to the power $1/p$.  Because $X$
is compact, the infimum is attained and $W_p(\mu,\nu) < \infty$
for all $\mu, \nu \in \calP(X)$.

\subsection{Metric Properties}

$W_p$ is a metric on $\calP(X)$:
\begin{itemize}
  \item \textbf{Non-negativity:} $W_p \geq 0$, with
    $W_p(\mu,\nu) = 0$ iff $\mu = \nu$.
  \item \textbf{Symmetry:} $W_p(\mu,\nu) = W_p(\nu,\mu)$,
    because coupling is symmetric.
  \item \textbf{Triangle inequality:} $W_p(\mu,\nu) \leq
    W_p(\mu,\sigma) + W_p(\sigma,\nu)$ for all $\sigma$.
    This follows from the \emph{gluing lemma}: given optimal
    couplings $\pi_1 \in \Pi(\mu,\sigma)$ and
    $\pi_2 \in \Pi(\sigma,\nu)$, one can construct a joint
    distribution on $X^3$ whose projections on the relevant
    pairs yield a coupling in $\Pi(\mu,\nu)$.
\end{itemize}
The metric $W_p$ metrises weak convergence of probability
measures on $X$ (when $X$ is compact, this is equivalent to
convergence in distribution).

\subsection{Ordering Across \texorpdfstring{$p$}{p}}

By Jensen's inequality applied to the concave function
$t \mapsto t^{p/q}$ for $p \leq q$:
\begin{equation}\label{eq:ordering}
  W_p(\mu,\nu) \leq W_q(\mu,\nu), \quad 1 \leq p \leq q.
\end{equation}
Thus $W_1 \leq W_2 \leq W_\infty$ and smaller $p$ gives a
weaker topology.

\subsection{The Case \texorpdfstring{$p = 1$}{p=1}:
            Kantorovich-Rubinstein}

For $p = 1$ the Kantorovich dual \eqref{eq:kantorovich_dual}
with $c = d$ simplifies to the celebrated
\emph{Kantorovich-Rubinstein} representation:
\begin{equation}\label{eq:w1_dual}
  \boxed{W_1(\mu,\nu)
  = \sup_{\|h\|_{\mathrm{Lip}} \leq 1}
  \left[\int_X h\,d\mu - \int_X h\,d\nu\right],}
\end{equation}
where the supremum is over all $1$-Lipschitz functions
$h: X \to \R$ (i.e., $|h(x) - h(y)| \leq d(x,y)$).
$W_1$ is also called the \emph{earth mover's distance}
because it measures the minimum work needed to redistribute
a pile of earth from shape $\mu$ to shape $\nu$.

\subsection{The Case \texorpdfstring{$p = 2$}{p=2}:
            Brenier Maps and Displacement Interpolation}

For $p = 2$ on $\R^d$ (or on a Riemannian manifold), the
optimal coupling under the quadratic cost is, under mild
regularity, a deterministic coupling induced by the
\emph{Brenier map} $T = \nabla\phi$ (gradient of a convex
potential $\phi$): the unique optimal map $T$ such that
$T_\#\mu = \nu$ \citep{Brenier1991}.

A remarkable feature of $W_2$ is that it supports a
Riemannian-like geometry on $\calP(X)$.  The
\emph{displacement interpolation} \citep{McCann1997}
$\mu_t = ((1-t)\,\mathrm{id} + t\,T)_\#\mu$ is the geodesic
(shortest path) between $\mu$ and $\nu$ in the $W_2$ metric.
This provides a natural notion of ``straight-line
interpolation'' between probability distributions, with no
analogue in the KL framework.

The $W_2$ metric also arises naturally in the
\emph{Jordan-Kinderlehrer-Otto (JKO) scheme}
\citep{Jordan1998}: the Fokker-Planck equation for a diffusion
process is the gradient flow of the free energy functional
$\mu \mapsto \KL{\mu}{\mu_\infty}$ (where $\mu_\infty$ is the
invariant measure) with respect to the $W_2$ metric.  This
deep connection between gradient flows, KL divergence, and $W_2$
is the subject of \citet{AmbrosioGigliSavare2008}.

\subsection{Summary of Properties}

Table~\ref{tab:properties} collects key properties for
comparison.

\begin{table}[ht]
\centering
\caption{Properties of $\KL{f}{g}$ and $W_p(\mu,\nu)$, where $f,g$
are mutually absolutely continuous densities on a compact metric space
$(X,d)$.}
\label{tab:properties}
\medskip
\begin{tabular}{p{4.5cm}ll}
\toprule
Property & $\KL{f}{g}$ & $W_p(\mu,\nu)$ \\
\midrule
Non-negative & Yes & Yes \\
Symmetric & No & Yes \\
Triangle inequality & No & Yes \\
Metric & No (divergence only) & Yes \\
Requires $f \ll g$ for finiteness & Yes & No \\
Sensitive to ground metric $d$ on $X$ & No & Yes \\
Induces geometry on $\calP(X)$ & Yes (Fisher--Rao) & Yes (Otto) \\
Defined via coupling & No & Yes \\
Dual representation & Donsker--Varadhan & Kantorovich--Rubinstein \\
Convex in $(f,g)$ jointly & Yes & Yes (convex in each marginal) \\
Gradient-flow structure & KL = functional in $W_2$-GF & $W_2$ = metric \\
\bottomrule
\end{tabular}
\end{table}

%======================================================================
\section{Comparing Kullback-Leibler and Wasserstein}
\label{sec:comparison}

\subsection{Sensitivity to the Ground Metric on $X$}

A key difference between KL divergence and Wasserstein distance
is their relationship to the \emph{ground metric} $d$ on the
sample space $X$---to be distinguished carefully from the
geometry each object induces on the \emph{space of probability
distributions} $\calP(X)$ (which we discuss in
Section~\ref{sec:info_geom}).

The Wasserstein distance $W_p(\mu,\nu)$ is constructed from the
ground metric $d$ on $X$: moving mass from $x$ to $y$ costs
$d(x,y)^p$.  Consequently, if $\mu = \delta_{x_0}$ and
$\nu = \delta_{y_0}$ are point masses, then
$W_p(\mu,\nu) = d(x_0,y_0)^p > 0$ whenever $x_0 \neq y_0$,
and $W_p \to 0$ as $d(x_0, y_0) \to 0$.  The Wasserstein
distance encodes the ground metric.

The KL divergence does not use $d$ at all.  In the same example,
$\mu = \delta_{x_0}$ and $\nu = \delta_{y_0}$ are not absolutely
continuous with respect to each other ($\delta_{x_0} \not\ll
\delta_{y_0}$), so $\KL{\mu}{\nu} = +\infty$ regardless of how
small $d(x_0,y_0)$ is.

As a continuous example, consider on $X = [0,1]$ the densities
$f_n(x) = 1 + \sin(2\pi n x)$ and $g_n(x) = 1 - \sin(2\pi n x)$.
Both are valid probability densities for all $n \geq 1$.  As
$n \to \infty$, the two oscillate faster and faster; they become
equally spread over $[0,1]$ but always ``out of phase.''  It can
be shown that $W_1(f_n, g_n) \to 0$ (the mass only needs to
travel $\approx 1/n$ to cancel the oscillations), while
$\KL{f_n}{g_n}$ remains bounded away from zero (the pointwise
ratio $f_n/g_n$ continues to oscillate significantly).

\subsection{Information Geometry: The Geometry KL Induces on $\calP(X)$}
\label{sec:info_geom}

Although KL divergence ignores the ground metric $d$ on $X$,
it is far from geometry-free: it \emph{generates} a Riemannian
geometry on the statistical manifold $\calP(X)$ through its
infinitesimal expansion.  This is the subject of
\emph{information geometry} \citep{Amari1985,AmariNagaoka2000,Amari2016,Nielsen2020}.

Let $\mathcal{M} = \{p(\cdot;\theta) : \theta \in \Theta\}$ be
a smooth parametric family of probability densities.  A
straightforward Taylor expansion of the KL divergence in a
small perturbation $\theta \to \theta + d\theta$ gives:
\begin{equation}\label{eq:fisher_kl}
  \KL{p(\cdot;\theta)}{p(\cdot;\theta+d\theta)}
  = \tfrac{1}{2}\,g_{ij}(\theta)\,d\theta^i d\theta^j + O(|d\theta|^3),
\end{equation}
where the coefficients
\begin{equation}\label{eq:fisher_metric}
  g_{ij}(\theta)
  = \mathbb{E}_{\theta}\!\left[
    \frac{\partial \log p(X;\theta)}{\partial \theta_i}
    \frac{\partial \log p(X;\theta)}{\partial \theta_j}
  \right]
  = -\mathbb{E}_{\theta}\!\left[
    \frac{\partial^2 \log p(X;\theta)}{\partial \theta_i \partial \theta_j}
  \right]
\end{equation}
form the \emph{Fisher information matrix}.  The matrix
$(g_{ij}(\theta))$ is positive semidefinite (positive definite
under identifiability conditions) and defines a Riemannian metric
on $\mathcal{M}$---the \emph{Fisher-Rao metric}
\citep{Rao1945}.  In this sense, the KL divergence is the
\emph{potential} whose Hessian generates the geometry of the
statistical manifold.

Amari \citep{Amari1985,AmariNagaoka2000} enriched this Riemannian
structure with a family of \emph{$\alpha$-connections} $\nabla^{(\alpha)}$
for $\alpha \in [-1,1]$.  The $e$-connection ($\alpha = +1$,
exponential) and the $m$-connection ($\alpha = -1$, mixture) are
dually flat with respect to the Fisher metric.  These connections
encode the natural coordinate systems for exponential and mixture
families respectively---and the KL divergence is the
\emph{divergence function} that generates the dual structure via:
\[
  \KL{p_0}{p_1}
  = \tfrac{1}{2}\,g_{ij}\,d\theta^i d\theta^j
    - \tfrac{1}{3}T_{ijk}\,d\theta^i d\theta^j d\theta^k + \cdots,
\]
where $T_{ijk}$ (the \emph{Amari-Chentsov tensor}) encodes the
skewness of the deviation from Riemannian symmetry.

A central operation in information geometry is the
\emph{$m$-projection} (mixture projection) and
\emph{$e$-projection} (exponential projection):
\[
  \Pi_e(p; \mathcal{F}) = \arg\min_{q \in \mathcal{F}} \KL{p}{q},
  \qquad
  \Pi_m(p; \mathcal{F}) = \arg\min_{q \in \mathcal{F}} \KL{q}{p},
\]
for a sub-manifold $\mathcal{F}$.  The Pythagorean theorem of
information geometry states that for dually flat spaces and a
convex sub-manifold $\mathcal{F}$:
\[
  \KL{p}{q} = \KL{p}{\Pi_e(p;\mathcal{F})} + \KL{\Pi_e(p;\mathcal{F})}{q},
  \qquad \forall\, q \in \mathcal{F},
\]
an exact analogue of the Euclidean Pythagorean theorem.
Frank Nielsen and collaborators \citep{Nielsen2020,Nielsen2010}
have developed efficient computational algorithms exploiting this
geometry, including Bregman-divergence-based Voronoi diagrams and
clustering in the KL geometry.

\paragraph{Two distinct geometries on $\calP(X)$.}
The Fisher-Rao geometry and the Wasserstein-2 (Otto) geometry are
\emph{both} valid Riemannian geometries on the space of probability
distributions, but they capture different phenomena:
\begin{itemize}
  \item The \textbf{Fisher-Rao geometry} is intrinsic to the
    \emph{distributional family}: it measures infinitesimal
    distinguishability between neighbouring distributions,
    independently of any metric structure on the sample space $X$.
    Geodesics in the Fisher-Rao metric are statistically natural
    (e.g., for exponential families they are straight lines in
    natural parameters).
  \item The \textbf{Otto (Wasserstein-2) geometry} is built from
    the ground metric $d$ on $X$: it measures the cost of
    transporting mass between distributions.  Geodesics in the
    $W_2$ metric are displacement interpolations (Section~\ref{sec:wasserstein}).
\end{itemize}
The two geometries are genuinely different: the Fisher-Rao geodesic
between two Gaussians passes through other Gaussians, while the
$W_2$ geodesic (displacement interpolation) generates intermediate
distributions that are also Gaussian but follow a different path
through parameter space \citep{Villani2009}.

\subsection{Absolute Continuity}

$\KL{f}{g}$ requires $\mu \ll \nu$ (i.e., $f \ll g$) to be
finite.  If $\mu$ assigns positive mass to a set that $\nu$
assigns zero mass, then $\KL{\mu}{\nu} = +\infty$.

$W_p(\mu,\nu)$ is defined and finite for all
$\mu, \nu \in \calP(X)$ when $X$ is compact, with no
absolute-continuity requirement.  The optimal coupling
$\pi^*$ may or may not have a density.

This distinction has important practical consequences.  When
learning a generative model, if the model distribution $\mu_\theta$
has support on a lower-dimensional manifold (e.g., a neural
network maps Gaussian noise to images), then $\KL{\mu_\theta}{\nu}$
may be $+\infty$ for most $\theta$, giving no useful gradient.
The Wasserstein distance is well-defined in these settings and
provides useful gradient information---the key observation behind
the Wasserstein GAN \citep[Ch.\ 7]{PeyreCuturi2019}.

\subsection{Dual Representations Compared}

The dual representations of $W_1$ and $\KL{f}{g}$ are worth
placing side by side:
\begin{align}
  W_1(\mu,\nu)
  &= \sup_{\|h\|_{\mathrm{Lip}} \leq 1}
    \Bigl[\E{\mu}{h} - \E{\nu}{h}\Bigr],
  \label{eq:w1_dual2}\\[6pt]
  \KL{f}{g}
  &= \sup_{h \in L^1(\mu)}
    \Bigl[\E{\mu}{h} - \log\E{\nu}{e^h}\Bigr].
  \label{eq:kl_dual2}
\end{align}
Both involve the difference of expectations of a test function
$h$ under the two measures.  The differences are:
\begin{itemize}
  \item $W_1$ uses a \emph{linear} test-function difference
    $\E{\mu}{h} - \E{\nu}{h}$, constrained to Lipschitz
    functions.  The constraint encodes the geometry of $X$
    through the Lipschitz condition.
  \item KL uses a \emph{non-linear} expression
    $\E{\mu}{h} - \log\E{\nu}{e^h}$.  The exponential
    transforms $h$ by an exponential twist of the measure
    $\nu$ (a Cram\'er tilt), which is sensitive to large
    deviations in the tails.  There is no geometric constraint
    on $h$.
\end{itemize}
The linear structure of $W_1$'s dual makes it well-behaved
for stochastic gradient methods and stable under small geometry-
preserving perturbations.  The exponential structure of KL's
dual makes it powerful for hypothesis testing and large-deviations
analysis, but unstable when the measures have very different
tails.

\subsection{The (Absence of) Coupling Underneath KL Divergence}
\label{sec:kl_coupling}

Recall that the Wasserstein distance is defined as an infimum over the set $\Pi(\mu,\nu)$ of all couplings of $\mu$ and $\nu$.

Does $\KL{f}{g}$ come from a similar optimisation over
$\Pi(\mu,\nu)$?  In general, \textbf{no}.  The KL divergence
does not use a coupling of $f$ and $g$ at all.  Its definition
\eqref{eq:kl_def} requires only the two densities evaluated at
the \emph{same point} $x \in X$:
\[
  \KL{f}{g} = \int_X f(x)\log\frac{f(x)}{g(x)}\,d\lambda(x)
  = \E{\mu}{\log\frac{f(X)}{g(X)}}.
\]
There is only one random variable here, $X \sim \mu = f\lambda$;
no second variable $Y \sim \nu$ appears.  The comparison between
$f$ and $g$ is \emph{pointwise}: both densities are evaluated at
the same location $x$ using the common reference measure
$\lambda$, and their ratio $f(x)/g(x)$ is the log-likelihood
ratio.

In informal terms: KL divergence compares the two mountains $f$
and $g$ by looking at their heights $f(x)$ and $g(x)$ at each
location $x$.  Wasserstein compares them by
asking how much effort it takes to rearrange the earth in one
mountain to create the other.  KL cares about altitudes;
Wasserstein cares about lateral displacement.

One way to see why no coupling naturally underlies KL: a coupling
$\pi \in \Pi(\mu,\nu)$ requires the two marginals to be the
\emph{respective} $x$- and $y$-marginals.  In KL, there is only
one variable: both $f$ and $g$ are functions of the same
variable $x$.  The ``natural'' candidate for an underlying
coupling---the \emph{diagonal coupling} $\pi_\Delta(x,y) =
f(x)\delta(y-x)$---has marginals $(\mu, \mu)$, not $(\mu, \nu)$,
and so is not in $\Pi(\mu,\nu)$ unless $\mu = \nu$.

One can formally compute the KL divergence of the diagonal
coupling from the independent coupling:
\begin{align}\label{eq:kl_diagonal}
  \KL{\pi_\Delta}{\mu\otimes\nu}
  &= \int_{X\times X} \pi_\Delta(x,y)\log
    \frac{\pi_\Delta(x,y)}{f(x)g(y)}\,d\lambda(x)\,d\lambda(y)
  \nonumber\\
  &= \int_X f(x)\log\frac{f(x)}{f(x)g(x)}\,d\lambda(x)
  = -\int_X f(x)\log g(x)\,d\lambda(x)
  = H_{\mathrm{cross}}(f,g),
\end{align}
which is the \emph{cross-entropy}, not $\KL{f}{g}$.  The
difference between cross-entropy and KL is precisely the entropy:
$H_{\mathrm{cross}}(f,g) = H(f) + \KL{f}{g}$.  So the diagonal
coupling contributes the entropy term $H(f)$ as overhead, along
with $\KL{f}{g}$.  This shows that even if one forces a
coupling interpretation, there is an irreducible ``entropic cost''
$H(f)$ that has nothing to do with the discrepancy between $f$
and $g$.

\paragraph{Summary.}
The coupling perspective is fundamental to Wasserstein:
$W_p = \inf_\pi [\cdots]$ over $\Pi(\mu,\nu)$.  KL divergence
has no such coupling: it is a functional of the ratio $f/g$ at
the same point, using the common reference measure $\lambda$.
The two objects are measuring fundamentally different things, and
the question of ``which coupling underlies KL'' does not have
a natural answer in the optimal transport sense.

\subsection{Entropic Regularisation: KL as a Penalty on Couplings}
\label{sec:entropic_ot}

Although KL divergence does not itself arise from optimising
over $\Pi(\mu,\nu)$, it plays a central role as a
\emph{regulariser} when added to the Kantorovich objective.
The \emph{entropy-regularised optimal transport} problem,
popularised by \citet{Cuturi2013}, is:
\begin{equation}\label{eq:entropic_ot}
  \boxed{\mathrm{OT}_\varepsilon(\mu,\nu)
  = \inf_{\pi \in \Pi(\mu,\nu)}
  \left[\int_{X\times X} c(x,y)\,d\pi(x,y)
  + \varepsilon\,\KL{\pi}{\mu\otimes\nu}\right],
  \quad \varepsilon > 0.}
\end{equation}
Here $\KL{\pi}{\mu\otimes\nu}$ is the KL divergence of the
joint $\pi$ from the \emph{independent} coupling $\mu\otimes\nu$:
\begin{equation}\label{eq:kl_joint}
  \KL{\pi}{\mu\otimes\nu}
  = \int_{X\times X} \pi(x,y)\log
    \frac{\pi(x,y)}{f(x)g(y)}\,d\lambda(x)\,d\lambda(y)
  = I_\pi(X;Y),
\end{equation}
the \emph{mutual information} between $X$ and $Y$ under the
coupling $\pi$.  The penalty $\varepsilon\, I_\pi(X;Y)$
discourages couplings that create strong statistical dependence
between $X$ and $Y$; the independent coupling $\mu\otimes\nu$
achieves $I = 0$ and incurs no penalty.

The solution to \eqref{eq:entropic_ot} is a coupling
$\pi_\varepsilon^* \in \Pi(\mu,\nu)$ with density
\begin{equation}\label{eq:sinkhorn_kernel}
  \pi_\varepsilon^*(x,y) \propto
  \exp\!\left(-\frac{c(x,y)}{\varepsilon}\right) f(x)g(y),
\end{equation}
up to normalisation functions $u(x)$ and $v(y)$ (Sinkhorn
scaling).  The \textsc{Sinkhorn} algorithm
\citep{Cuturi2013} computes $u$ and $v$ by alternating
projections, converging geometrically.

\paragraph{Two limits of the regularisation parameter.}
\begin{enumerate}
  \item \textbf{As $\varepsilon \to 0$}: $\pi_\varepsilon^*$
    concentrates on the optimal transport plan of the
    unregularised problem \eqref{eq:kantorovich}, and
    $\mathrm{OT}_\varepsilon \to C(\mu,\nu;c)$ (the unregularised
    OT cost, i.e., $W_p^p$ when $c = d^p$).  The penalty term
    vanishes, and the entropy of $\pi_\varepsilon^*$ tends to
    the (possibly small) entropy of the optimal transport plan.
  \item \textbf{As $\varepsilon \to \infty$}: the transport-cost
    term becomes negligible relative to the KL penalty, and
    $\pi_\varepsilon^* \to \mu\otimes\nu$ (the independent
    coupling).  Intuitively, when the penalty for statistical
    dependence is very large, the only affordable coupling is the
    independent one, regardless of how much transport cost it saves.
    The problem is dominated by entropy maximisation subject to the
    marginal constraints, and the solution is the independent coupling.
\end{enumerate}

\paragraph{Sinkhorn divergence.}
The raw entropy-regularised cost $\mathrm{OT}_\varepsilon(\mu,\nu)$
is neither symmetric nor zero when $\mu = \nu$.  The
\emph{Sinkhorn divergence} \citep{GenevayEtal2018} corrects this:
\begin{equation}\label{eq:sinkhorn_div}
  S_\varepsilon(\mu,\nu)
  = \mathrm{OT}_\varepsilon(\mu,\nu)
  - \tfrac{1}{2}\mathrm{OT}_\varepsilon(\mu,\mu)
  - \tfrac{1}{2}\mathrm{OT}_\varepsilon(\nu,\nu).
\end{equation}
$S_\varepsilon$ is symmetric, non-negative, zero iff $\mu = \nu$,
and interpolates between $W_p$ (as $\varepsilon \to 0$) and the
squared maximum mean discrepancy (MMD) (as $\varepsilon \to \infty$).
It enjoys the computational efficiency of entropy-regularised OT
while retaining the metric flavour of Wasserstein.

\subsection{The Schr\"odinger Bridge Problem}
\label{sec:schrod}

A related but distinct connection between couplings and KL
divergence is the \emph{Schr\"odinger bridge} problem, introduced
by \citet{Schrodinger1931} in the context of statistical mechanics.

Let $R \in \calP(X \times X)$ be a \emph{reference} joint
distribution (not necessarily in $\Pi(\mu,\nu)$) that encodes
some prior knowledge about how particles might pair up---for
example, the joint distribution obtained from a Brownian bridge
between two time points.  The Schr\"odinger bridge problem asks:
\begin{equation}\label{eq:schrod}
  \pi^* = \arg\min_{\pi \in \Pi(\mu,\nu)} \KL{\pi}{R},
\end{equation}
i.e., find the coupling that is closest to $R$ in KL divergence
while having the prescribed marginals $\mu$ and $\nu$.  This is a
projection of $R$ onto the marginal-constrained set $\Pi(\mu,\nu)$
in the KL sense.

The Schr\"odinger bridge is notable for several reasons:
\begin{enumerate}
  \item \textbf{KL on joints.}  Unlike $W_p$, which minimises
    a transport cost $\int c\,d\pi$, the Schr\"odinger bridge
    minimises KL divergence between joint distributions.  KL
    now appears as the objective, not a regulariser.
  \item \textbf{Connection to entropic OT.}  When
    $R(dx,dy) = e^{-c(x,y)/\varepsilon}\mu(dx)\nu(dy)/Z$
    (i.e., when $R$ is the law of a pair with density
    proportional to the transport kernel), the Schr\"odinger
    bridge $\pi^*$ coincides with the entropy-regularised OT
    solution $\pi^*_\varepsilon$ in \eqref{eq:sinkhorn_kernel}.
  \item \textbf{Convergence to Monge.}  When $R$ is the
    law of $(X_0, X_1)$ under a Brownian motion with diffusion
    coefficient $\sigma > 0$, as $\sigma \to 0$ the
    Schr\"odinger bridge converges to the Monge optimal
    transport plan.
\end{enumerate}
A comprehensive survey of the Schr\"odinger bridge and its
connections to optimal transport is given by \citet{Leonard2014}.

\subsection{Inequalities Relating KL and Wasserstein}
\label{sec:inequalities}

Even though KL divergence and Wasserstein distance measure
different things, they can be bounded in terms of each other
through intermediate quantities, principally the total variation
distance.

\paragraph{Pinsker's inequality.}

The \emph{total variation distance} between $\mu$ and $\nu$ is:
\begin{equation}\label{eq:tv}
  \TV{\mu}{\nu} = \sup_{A \in \mathcal{B}(X)} |\mu(A) - \nu(A)|
  = \frac{1}{2}\int_X |f(x) - g(x)|\,d\lambda(x).
\end{equation}
Pinsker's inequality (see \citealt{CoverThomas2006}) gives:
\begin{equation}\label{eq:pinsker}
  \boxed{\TV{\mu}{\nu}^2 \leq \tfrac{1}{2}\,\KL{f}{g}.}
\end{equation}
The reverse does not hold without additional conditions: KL can
be large even when total variation is close to 1 (for distributions
that barely overlap).

On a compact space with diameter $D = \sup_{x,y\in X}d(x,y)$,
the Kantorovich-Rubinstein theorem and the bound
$\int h(f-g) \leq \|h\|_\infty\|f-g\|_{L^1}$ give:
\begin{equation}\label{eq:w1_tv}
  W_1(\mu,\nu) \leq D\,\TV{\mu}{\nu}.
\end{equation}
Combining \eqref{eq:pinsker} and \eqref{eq:w1_tv}:
\begin{equation}\label{eq:w1_kl}
  W_1(\mu,\nu) \leq D\,\sqrt{\tfrac{1}{2}\,\KL{f}{g}}.
\end{equation}
Thus on bounded spaces, Wasserstein-1 is controlled by the square
root of KL.  The chain of inequalities is:
\[
  W_1 \;\leq\; D\,\TV{\mu}{\nu} \;\leq\; D\sqrt{\tfrac{1}{2}\KL{f}{g}}.
\]

\paragraph{Talagrand's transportation inequality.}

The inequality \eqref{eq:w1_kl} goes in the direction
KL $\Rightarrow$ $W_1$.  For $W_2$, a related result holds
under a structural condition on one of the measures.

A measure $\nu$ is said to satisfy a \emph{$T_2$ transportation
inequality} with constant $C > 0$ if for all $\mu \in \calP(X)$:
\begin{equation}\label{eq:t2}
  W_2(\mu,\nu)^2 \leq C\,\KL{\mu}{\nu}.
\end{equation}
\citet{Talagrand1996} proved \eqref{eq:t2} for $\nu$ the
standard Gaussian on $\R^d$ with $C = 2$.  The key insight is
that the Gaussian satisfies a logarithmic Sobolev inequality
(LSI) with constant $\alpha = 1$, and LSI implies $T_2$ with
$C = 2/\alpha$.

The precise connection between LSI and $T_2$ was established by
\citet{OttoVillani2000} through the \emph{HWI inequality}:
\begin{equation}\label{eq:hwi}
  \KL{\mu}{\nu}
  \leq W_2(\mu,\nu)\sqrt{I(\mu|\nu)}
  - \frac{\alpha}{2}W_2(\mu,\nu)^2,
\end{equation}
where $I(\mu|\nu) = \int f|\nabla\log(f/g)|^2\,d\lambda$ is the
\emph{relative Fisher information}, and $\alpha$ is the
log-Sobolev constant of $\nu$.  The inequality \eqref{eq:hwi}
is called HWI because it relates $H = \KL{\mu}{\nu}$
(entropy), $W = W_2$ (Wasserstein),  and $I = I(\mu|\nu)$
(Fisher information).  Combining the HWI inequality with
Young's inequality yields Talagrand's $T_2$ from
\eqref{eq:t2}.

Thus under a log-Sobolev condition on $\nu$:
\[
  W_2(\mu,\nu)^2 \leq \frac{2}{\alpha}\,\KL{\mu}{\nu}.
\]
In the other direction, without further conditions, no inequality
of the form $\KL{\mu}{\nu} \lesssim W_2(\mu,\nu)^\alpha$ holds
globally on a compact space ($D_{\mathrm{KL}}$ can blow up while $W_2$
remains bounded).

Table~\ref{tab:inequalities} summarises the inequalities in
one place.

\begin{table}[ht]
\centering
\caption{Inequalities relating KL divergence, total variation,
and Wasserstein distance.  Here $D = \mathrm{diam}(X)$ and
$\alpha$ is the log-Sobolev constant of $\nu$.}
\label{tab:inequalities}
\medskip
\begin{tabular}{p{5.2cm}p{5.2cm}l}
\toprule
Inequality & Content & Reference \\
\midrule
Pinsker & $\TV{\mu}{\nu}^2 \leq \frac{1}{2}\KL{f}{g}$ &
  \citealt{CoverThomas2006} \\[4pt]
$W_1$--TV & $W_1(\mu,\nu) \leq D\,\TV{\mu}{\nu}$ &
  Kantorovich-Rubinstein \\[4pt]
$W_1$--KL & $W_1(\mu,\nu) \leq D\sqrt{\frac{1}{2}\KL{f}{g}}$ &
  Pinsker + KR \\[4pt]
Talagrand $T_2$ & $W_2(\mu,\nu)^2 \leq \frac{2}{\alpha}\KL{\mu}{\nu}$
  & \citealt{Talagrand1996} \\[4pt]
HWI & $\KL{\mu}{\nu} \leq W_2\sqrt{I(\mu|\nu)} -
  \frac{\alpha}{2}W_2^2$ & \citealt{OttoVillani2000} \\
\bottomrule
\end{tabular}
\end{table}

\paragraph{Remark on the sharpness of Pinsker.}
Pinsker's inequality is not tight in general.  For Gaussian
distributions $\mu = \mathcal{N}(m_1, \sigma^2)$ and
$\nu = \mathcal{N}(m_2, \sigma^2)$ with $\delta =
|m_1 - m_2|/(2\sigma)$ (half the normalised mean shift):
$\KL{\mu}{\nu} = \delta^2 / 2$ and
$\TV{\mu}{\nu} = 2\Phi(\delta) - 1 \approx 2\phi(0)\delta$
for small $\delta$.  So for small perturbations,
$\|f-g\|_{\mathrm{TV}} \approx \frac{2}{\sqrt{2\pi}}\sqrt{2\,\KL{\mu}{\nu}}$, while Pinsker
gives $\|f-g\|_{\mathrm{TV}} \leq \sqrt{\KL{\mu}{\nu}/2}$.  The two are of the same order
but differ by a constant.

%======================================================================
\section{Conclusions}

We summarise the main findings:

\begin{enumerate}
  \item \textbf{KL divergence is insensitive to the ground metric on $X$,
    but induces its own geometry on $\calP(X)$.}
    As an $f$-divergence, KL compares the two densities at the
    \emph{same point} $x$ via the ratio $f(x)/g(x)$, using the
    common reference measure $\lambda$ as a shared coordinate system.
    No ground-metric cost $d(x,y)$ enters.  However, KL is far
    from geometry-free: its Hessian generates the Fisher information
    metric on any statistical manifold
    $\{p(\cdot;\theta)\}$, the foundational object of Amari's
    information geometry \citep{Amari2016,AmariNagaoka2000}.  This
    is an intrinsic geometry on the space of distributions,
    distinct from the ground-metric-derived $W_2$ (Otto) geometry.
    No coupling of $\mu$ and $\nu$ is involved in KL.

  \item \textbf{Wasserstein distance is defined by optimising
    over couplings.}  Every element of $\Pi(\mu,\nu)$ is a
    feasible coupling; the Wasserstein distance picks the
    cheapest one.  The optimal coupling encodes the
    geometry-respecting ``best matching'' of mass in $\mu$ to
    mass in $\nu$.

  \item \textbf{No coupling naturally underlies KL.}  The
    ``diagonal coupling'' $\pi_\Delta(x,y) = f(x)\delta(y-x)$
    has marginals $(\mu,\mu)$ rather than $(\mu,\nu)$, and so
    lies in $\Pi(\mu,\mu)$ rather than $\Pi(\mu,\nu)$.  Its
    KL divergence from the product measure equals the
    cross-entropy $H_{\mathrm{cross}}(f,g) = H(f) + \KL{f}{g}$,
    not $\KL{f}{g}$ itself.

  \item \textbf{KL emerges as a regulariser on couplings.}
    In entropy-regularised OT \eqref{eq:entropic_ot}, KL
    divergence of the joint $\pi$ from the independent coupling
    $\mu\otimes\nu$ appears as a penalty.  This penalises
    couplings that create statistical dependence between the
    two marginals, and its role is to smooth the discrete
    optimal transport plan into a continuous kernel.

  \item \textbf{KL emerges as an objective on joints in the
    Schr\"odinger bridge.}  The Schr\"odinger bridge
    \eqref{eq:schrod} minimises KL divergence of a coupling
    $\pi$ from a reference joint $R$, subject to marginal
    constraints.  This is the natural dual to OT: where OT
    fixes a transport cost and optimises freely over $\Pi$,
    the Schr\"odinger bridge fixes a reference joint and
    optimises according to an entropic (KL) criterion.

  \item \textbf{Inequalities connect the two.}  On compact
    spaces, $W_1 \leq D\sqrt{D_{\mathrm{KL}}/2}$ (via Pinsker and
    Kantorovich-Rubinstein), and $W_2^2 \leq (2/\alpha)D_{\mathrm{KL}}$
    (Talagrand) when $\nu$ satisfies a log-Sobolev inequality.
    These bounds go in the direction KL $\Rightarrow$
    Wasserstein; the reverse typically fails without additional
    conditions.

  \item \textbf{Both are special cases of a richer hierarchy.}
    The $f$-divergence family (of which KL is a member) and the
    class of Integral Probability Metrics (of which $W_1$ is a
    member) do not overlap.  Their interaction is mediated by
    total variation distance, which belongs to both families.
    The entropy-regularised OT and the Schr\"odinger bridge
    occupy an intermediate position, blending variational
    optimisation over couplings with entropic functionals.
\end{enumerate}

%======================================================================
\appendix

\section{KL Divergence and Mutual Information}
\label{app:mi}

Given a coupling $\pi \in \Pi(\mu,\nu)$ with joint density
$\pi(x,y)$, the \emph{mutual information} between $X$ and $Y$
under $\pi$ is:
\begin{equation}\label{eq:mi}
  I_\pi(X;Y) = \KL{\pi}{\mu\otimes\nu}
  = \int_{X\times X} \pi(x,y)\log
    \frac{\pi(x,y)}{f(x)g(y)}\,d\lambda(x)\,d\lambda(y).
\end{equation}
This is the KL divergence of the joint from the independent
coupling.  It measures how much statistical dependence the
coupling $\pi$ introduces beyond what the two marginals alone
would suggest.

Key properties of $I_\pi(X;Y)$:
\begin{itemize}
  \item $I_\pi(X;Y) \geq 0$, with equality iff $X \perp Y$
    under $\pi$ (i.e., $\pi = \mu \otimes \nu$).
  \item Among all couplings $\pi \in \Pi(\mu,\nu)$, the
    \emph{independent coupling} minimises $I_\pi(X;Y)$,
    achieving the value $0$.  It is the least statistically
    dependent coupling.
  \item The \emph{optimal transport coupling} $\pi^*_W$
    (for $W_p$) typically maximises statistical dependence
    for a given marginal pair: the Monge map $Y = T(X)$
    achieves $I_\pi(X;Y) = H(Y)$ (if $T$ is bijective and
    $\mu$ absolutely continuous), since $Y$ is fully determined
    by $X$ and $H(Y|X) = 0$.
\end{itemize}
The contrast is striking: the OT coupling maximises dependence
(all information about $X$ is encoded in $Y$), while the
independent coupling minimises it.  The entropic regulariser
$\varepsilon I_\pi$ in \eqref{eq:entropic_ot} penalises the
OT coupling for being ``too deterministic'' and drives the
solution toward the independent coupling as $\varepsilon$
increases.

\section{The Special Case of Gaussian Distributions}
\label{app:gaussian}

Let $\mu = \mathcal{N}(m_1, \Sigma_1)$ and
$\nu = \mathcal{N}(m_2, \Sigma_2)$ be Gaussian distributions
on $\R^d$.  Both the KL divergence and the Wasserstein-2
distance have closed forms.

\paragraph{KL divergence.}
\begin{equation}\label{eq:kl_gaussian}
  \KL{\mu}{\nu}
  = \tfrac{1}{2}\Bigl[
    \mathrm{tr}(\Sigma_2^{-1}\Sigma_1)
    + (m_2 - m_1)^\top \Sigma_2^{-1}(m_2 - m_1)
    - d
    + \log\frac{\det\Sigma_2}{\det\Sigma_1}
  \Bigr].
\end{equation}

\paragraph{Wasserstein-2 distance (Bures metric).}
The $W_2^2$ distance between Gaussians, sometimes called the
Bures(-Wasserstein) metric, is \citep{Villani2009}:
\begin{equation}\label{eq:w2_gaussian}
  W_2(\mu,\nu)^2
  = \|m_1 - m_2\|^2
  + \mathcal{B}(\Sigma_1, \Sigma_2)^2,
\end{equation}
where the Bures distance between covariances is
\begin{equation}\label{eq:bures}
  \mathcal{B}(\Sigma_1, \Sigma_2)^2
  = \mathrm{tr}\,\Sigma_1 + \mathrm{tr}\,\Sigma_2
  - 2\,\mathrm{tr}\Bigl[
    \bigl(\Sigma_1^{1/2}\Sigma_2\Sigma_1^{1/2}\bigr)^{1/2}
  \Bigr].
\end{equation}
For diagonal $\Sigma_1 = \sigma_1^2 I$ and
$\Sigma_2 = \sigma_2^2 I$: $\mathcal{B}^2 = d(\sigma_1 - \sigma_2)^2$
and $W_2(\mu,\nu)^2 = \|m_1-m_2\|^2 + d(\sigma_1-\sigma_2)^2$.

The Gaussian case illustrates the geometric sensitivity of
$W_2$: it depends linearly on the distance between means
($\|m_1-m_2\|$) and on the distance between standard
deviations ($|\sigma_1-\sigma_2|$) in the isotropic case.
The KL divergence \eqref{eq:kl_gaussian} is not symmetric
and grows quadratically in $\Sigma_2^{-1/2}(m_1-m_2)$,
favouring the direction of low variance in $\nu$.

\bigskip
\bibliographystyle{plainnat}
\bibliography{ot_kl_divergence}

\end{document}
